% ============================================================================
% 06_arquitectura_vistas.tex — Cuerpo del entregable "Documento de
% Arquitectura de Software: Vistas"
% Sin preámbulo: se incluye tanto desde 06_Arquitectura_Vistas.tex
% (standalone) como, eventualmente, desde un documento maestro.
% ============================================================================
\section{Introducción}

\subsection{Propósito y alcance}

Este documento describe la arquitectura de software de \textbf{LSA Salud}
mediante el modelo de vistas \textbf{4+1} de Kruchten, aplicando la
metodología de diseño usada en el resto de los entregables (Larman / Proceso
Unificado, notación UML 2.5 --- ver Referencias de la ERS, Entregable 05).
Incluye los componentes que realizan los requisitos RF-03, RF-04 y
RF-06 (consentimiento, auditoría y reporte de traducción). Mientras que el Modelo de Dominio (Entregable 01) describe los conceptos del
\emph{problema}, este documento describe las decisiones de \emph{diseño} del
software que lo resuelve: módulos, procesos concurrentes, organización del
código fuente y despliegue físico.

Se presentan cuatro vistas independientes más una vista de escenarios que
las integra:
\begin{itemize}
    \item \textbf{Vista Lógica} (\ref{sec:vista-logica}): módulos de diseño
    del frontend y del backend, y sus dependencias.
    \item \textbf{Vista de Procesos} (\ref{sec:vista-procesos}): concurrencia
    en tiempo de ejecución entre el navegador, el backend y los procesos
    locales invocados.
    \item \textbf{Vista de Desarrollo} (\ref{sec:vista-desarrollo}):
    organización del repositorio y dependencias entre sus paquetes.
    \item \textbf{Vista Física} (\ref{sec:vista-fisica}): nodos de despliegue,
    procesos en ejecución, puertos y protocolos.
    \item \textbf{Vista de Escenarios} (\ref{sec:vista-escenarios}): traza
    CU-01 y CU-04 a través de las cuatro vistas anteriores.
\end{itemize}

Este documento no repite contenido de otros entregables: el detalle de la
persistencia basada en archivos se documenta en el Entregable 09; los IDs de
actores, casos de uso, RF y RNF se reutilizan tal como están definidos en
\texttt{IDS.md} y en la ERS (Entregable 05).

\subsection{Convenciones}

Los diagramas de este documento son de \textbf{diseño}, no de dominio: a
diferencia del Entregable 01, sí aparecen aquí \texttt{hooks} de React,
endpoints de FastAPI, procesos del sistema operativo y puertos de red, porque
son precisamente el objeto de estudio de una vista arquitectónica. Las
etiquetas de figuras y tablas de este entregable usan el prefijo
\texttt{arq-} para no colisionar con las de otros entregables.

% ============================================================================
\section{Vista Lógica}
\label{sec:vista-logica}

La Vista Lógica muestra los módulos de diseño de LSA Salud y sus
dependencias: en el frontend, las páginas de \texttt{app/}, los
\texttt{hooks} de React y los \texttt{components}; en el backend, el
módulo de endpoints (\texttt{app.py}), el modelo LSTM
(\texttt{src.model.lstm}) y el módulo de validación de capturas
(\texttt{capture\_validation}). El diálogo de consentimiento (P-05, RF-03.1)
no es un componente ni una ruta aparte: vive dentro de
\texttt{paciente/page} (\texttt{PacientePage}), que gatea el montaje de los
\texttt{hooks} de captura con la prop \texttt{enabled}. El registro de
auditoría (RF-04.1) tampoco es un módulo propio: son funciones
(\texttt{\_log\_audit}, \texttt{\_append\_jsonl}) dentro de \texttt{app.py}
mismo, invocadas por cada \textit{handler} que completa una operación de
sistema, y por el endpoint \texttt{POST /audit} para los eventos que origina
el frontend (\texttt{consent}, \texttt{translation\_displayed}). El reporte
de traducción (RF-06.1) se inicia desde \texttt{ChatBubble} y viaja por la
dependencia HTTP de \texttt{components} hacia \texttt{app.py}.

\diagrama{06_vista_logica}{\textwidth}{Vista Lógica de LSA Salud: módulos de diseño y dependencias}{fig:arq-vista-logica}

Puntos relevantes del diagrama:
\begin{itemize}
    \item Las páginas \texttt{paciente/page} (P-02) y \texttt{recepcion/page}
    (P-03) son los únicos consumidores directos de los \texttt{hooks} y de
    los \texttt{components}; \texttt{page} (P-01) solo navega hacia ellas.
    \item \texttt{useMediaPipe} extrae, en cada frame, los landmarks de
    manos, pose y rostro (\textit{blendshapes}) mediante MediaPipe, apoyado
    en las funciones puras de \texttt{lib/mediapipe.ts}. Los
    \textit{blendshapes} faciales aportan los marcadores no manuales de la
    LSA (preguntas, negación, intensidad) que forman parte de las
    características de reconocimiento (ver decisión arquitectónica en la
    sección~\ref{sec:decisiones}).
    \item Solo tres \texttt{hooks} cruzan la frontera hacia el backend:
    \texttt{useSignWebSocket} (WebSocket de predicción),
    \texttt{useSentenceBuilder} (traducción de oraciones) y
    \texttt{CaptureMode} (captura de muestras de entrenamiento y consulta de
    clases/estadísticas). El resto de los \texttt{hooks} y \texttt{components}
    es puramente de interfaz de usuario o de sincronización intra-navegador
    (\texttt{useChatChannel}, ver Vista de Procesos).
    \item En el backend, \texttt{app.py} concentra todos los endpoints y
    delega la clasificación en el módulo del modelo LSTM y la validación de
    capturas en \texttt{capture\_validation}, ambos desacoplados como
    módulos propios (RNF-T-03).
    \item El gateo de consentimiento (frontend, dentro de
    \texttt{paciente/page}) y las funciones de auditoría
    (\texttt{\_log\_audit} / \texttt{\_append\_jsonl}, backend, invocadas
    desde cada \textit{handler} de \texttt{app.py} que realiza una
    operación de sistema) siguen el mismo patrón de desacoplamiento que
    \texttt{capture\_validation}: ninguno de los \textit{handlers}
    necesita conocer el formato interno del registro de auditoría, solo
    invocar una función.
\end{itemize}

% ============================================================================
\section{Vista de Procesos}
\label{sec:vista-procesos}

La Vista de Procesos muestra la concurrencia en tiempo de ejecución del
sistema: el bucle de captura del navegador, la comunicación asincrónica por
WebSocket, el manejo diferenciado de solicitudes en el backend FastAPI, y la
sincronización entre la Estación del Paciente y la Estación de Recepción por
\texttt{BroadcastChannel}, que es completamente independiente del backend.

Se presenta un diagrama por proceso. Todos siguen la misma disposición: el
Paciente y la Estación del Paciente a la izquierda, los procesos del backend
y el canal de la estación en el centro, y la Estación de Recepción y el
Personal de Recepción o Enfermería a la derecha. En los procesos de gestión
del modelo, el Desarrollador de IA / Equipo de Proyecto se ubica a la
derecha. Cada actor y cada proceso aparece una sola vez por diagrama.

\clearpage
\subsection{Proceso 1: consentimiento de cámara}

Antes de activar la cámara, la Estación del Paciente muestra el diálogo de
consentimiento (CU-09, RF-03). La cámara solo se habilita si el paciente
acepta; en ambos casos la decisión se registra como evento anónimo de
auditoría.

\diagrama{06_proceso_1_consentimiento}{0.8\textwidth}{Proceso 1: consentimiento de cámara}{fig:arq-proceso-1}

\clearpage
\subsection{Proceso 2: captura y reconocimiento de señas}

El bucle de captura del navegador corre de forma continua mientras la cámara
está activa y envía la ventana de frames por WebSocket cada 8 frames. El
\textit{event loop} de FastAPI recibe la ventana y delega la inferencia del
modelo a un hilo propio, de modo que sigue atendiendo otros mensajes mientras
dura la predicción (CU-01, RF-01, RF-05).

\diagrama{06_proceso_2_prediccion}{\textwidth}{Proceso 2: captura y reconocimiento de señas}{fig:arq-proceso-2}

\clearpage
\subsection{Proceso 3: cierre de oración, traducción y voz}

Al detectar el fin de la oración, la Estación del Paciente envía las palabras
reconocidas al backend. El \textit{handler} se ejecuta en el
\textit{threadpool} de Starlette, por lo que las llamadas a Ollama y a Piper
TTS no bloquean el \textit{event loop}. Ollama solo se invoca cuando no hay
una plantilla médica ni una traducción en caché; si no está disponible, se
muestra la oración sin reconstruir (RNF-T-01), sin interrumpir el flujo
(CU-02, CU-03).

\diagrama{06_proceso_3_traduccion}{\textwidth}{Proceso 3: cierre de oración, traducción y voz}{fig:arq-proceso-3}

\clearpage
\subsection{Proceso 4: envío del mensaje del paciente al personal}

El mensaje traducido se propaga de la Estación del Paciente a la Estación de
Recepción por \texttt{BroadcastChannel}: un mecanismo del navegador entre
pestañas del mismo origen, que no pasa por el backend ni por ninguna red
(CU-01).

\diagrama{06_proceso_4_sincronizacion}{0.85\textwidth}{Proceso 4: envío del mensaje del paciente al personal}{fig:arq-proceso-4}

\clearpage
\subsection{Proceso 5: respuesta del personal}

El Personal de Recepción o Enfermería responde por escrito o por voz; el
dictado se transcribe en el propio navegador. La respuesta viaja por el mismo
canal hacia la Estación del Paciente (CU-04, CU-05, RF-02).

\diagrama{06_proceso_5_respuesta}{\textwidth}{Proceso 5: respuesta del personal}{fig:arq-proceso-5}

\clearpage
\subsection{Proceso 6: reporte de una traducción incorrecta}

El paciente o el personal pueden reportar un mensaje traducido desde su
propia estación. El backend valida el reporte, lo agrega a
\texttt{logs/reports.jsonl} y registra el evento de auditoría
correspondiente (CU-10, RF-06).

\diagrama{06_proceso_6_reporte}{\textwidth}{Proceso 6: reporte de una traducción incorrecta}{fig:arq-proceso-6}

\clearpage
\subsection{Proceso 7: fin de atención}

El Personal de Recepción o Enfermería finaliza la atención: su estación vacía
la conversación y envía el evento \texttt{CLEAR} por el canal, que vacía
también la Estación del Paciente. Ningún dato de la atención queda
persistido; solo se registra el evento anónimo de fin de atención (CU-12,
RNF-03).

\diagrama{06_proceso_7_fin_atencion}{\textwidth}{Proceso 7: fin de atención}{fig:arq-proceso-7}

\clearpage
\subsection{Proceso 8: captura de muestras de entrenamiento}

En el modo captura, el Desarrollador de IA / Equipo de Proyecto elige o crea
la seña y graba un conjunto de repeticiones. Cada muestra se valida en el
backend antes de guardarse en el sistema de archivos local (CU-06, CU-07).

\diagrama{06_proceso_8_captura_muestras}{0.85\textwidth}{Proceso 8: captura de muestras de entrenamiento}{fig:arq-proceso-8}

\clearpage
\subsection{Proceso 9: entrenamiento del modelo}

El entrenamiento corre como un proceso independiente del servicio, sin
afectar la atención en curso. Separa entrenamiento y validación por sesión de
grabación, guarda el modelo resultante y el servicio lo carga al reiniciarse
(CU-08).

\diagrama{06_proceso_9_entrenamiento}{0.85\textwidth}{Proceso 9: entrenamiento del modelo}{fig:arq-proceso-9}

% ============================================================================
\section{Vista de Desarrollo}
\label{sec:vista-desarrollo}

La Vista de Desarrollo mapea la organización física del código fuente en el
repositorio y las dependencias entre sus paquetes, relevante para la
compilación, el empaquetado y la asignación de responsabilidades de
mantenimiento (RNF-T-03).

\diagrama{06_vista_desarrollo}{\textwidth}{Vista de Desarrollo de LSA Salud: paquetes del repositorio y sus dependencias}{fig:arq-vista-desarrollo}

Puntos relevantes:
\begin{itemize}
    \item El frontend (\texttt{web/}) y el backend (\texttt{app.py},
    \texttt{src/}) son unidades de despliegue independientes, sin
    dependencia de compilación entre sí: se comunican únicamente en tiempo
    de ejecución, por HTTP y WebSocket.
    \item \texttt{train\_lstm.py} es un script de reentrenamiento
    \textbf{offline}, fuera del proceso del backend: lee \texttt{dataset/} y
    escribe en \texttt{models/}, sin invocación síncrona desde \texttt{app.py}
    (ver CU-08 en \texttt{IDS.md}).
    \item \texttt{eval\_model.py <directorio>} es otro script offline:
    importa el modelo ya cargado de \texttt{app.py} y evalúa su accuracy
    sobre un directorio de muestras \texttt{.npy}, sin exponer ningún
    endpoint ni modificar el proceso del backend.
    \item \texttt{dataset/}, \texttt{models/} y \texttt{voices/} son
    carpetas de datos, no paquetes de código: representan la persistencia
    basada en archivos del sistema (sin base de datos, ver
    Entregable 09 para el detalle del modelo de datos).
    \item \texttt{docs/} (esta documentación) no tiene dependencias de
    código y no participa en el flujo de ejecución.
    \item \texttt{logs/} es la carpeta de datos que recibe los eventos de
    auditoría (RF-04.1, \texttt{audit.jsonl}) y los reportes de traducción
    (RF-06.1, \texttt{reports.jsonl}), escritos por funciones dentro de
    \texttt{app.py} mismo, con el mismo patrón de persistencia basada en
    archivos que \texttt{dataset/} (ver Entregable 09).
\end{itemize}

% ============================================================================
\section{Vista Física}
\label{sec:vista-fisica}

La Vista Física muestra los nodos de despliegue, los procesos en ejecución
y los protocolos y puertos entre ellos, sobre la arquitectura de
\textbf{estaciones fijas múltiples} descripta en la ERS (2.1 y 3.3): el
sistema completo despliega tantas estaciones independientes como puntos de
atención tenga el sanatorio --- admisión, recepción, triaje, guardia ---,
cada una con su propio par de caras paciente--personal (doble pantalla)
sobre una misma PC. Cada estación es, por diseño, la materialización de las
dos caras de la estación fija de guardia/triaje exigida por RNF-04/RNF-07
(ver ERS, Sección~3.4.2); su ensamblaje físico sobre hardware
comercial y su calibración se realizan durante la instalación de cada
estación, fuera de alcance de software.

\diagrama{06_vista_fisica}{\textwidth}{Vista Física de LSA Salud: nodos, procesos y protocolos de despliegue de cada estación, replicados de forma independiente en admisión, recepción, triaje y guardia}{fig:arq-vista-fisica}

El diagrama detalla por completo una estación de referencia (Estación de
Admisión) y resume las demás (Recepción, Triaje, Guardia) como instancias
con la misma estructura interna, para mantener la figura legible: las
cuatro son, en tiempo de ejecución, despliegues idénticos y completamente
independientes entre sí.

Puntos relevantes:
\begin{itemize}
    \item Cada estación es una instancia local autocontenida: navegador,
    servidor Next.js (puerto 3000), backend FastAPI sobre Uvicorn (puerto
    8000) y Ollama (puerto 11434) corren en la misma PC de esa estación. No
    hay ningún nodo ni proceso externo a esa máquina, ni un servidor
    central compartido entre estaciones (RNF-T-02, RNF-03.1).
    \item Las estaciones no comparten datos entre sí: cada una conserva su
    propia conversación en memoria únicamente mientras dura la atención en
    curso, y la borra por completo al finalizarla (CU-12), sin persistirla
    en ningún momento (RNF-03). El backend de una estación restringe además
    los orígenes HTTP permitidos a la red local del sanatorio (RNF-T-06).
    \item Piper TTS \textbf{no} es un proceso separado con puerto propio: se
    carga como modelo ONNX embebido dentro del mismo proceso Python que
    corre \texttt{app.py}, a diferencia de Ollama, que sí es un servicio
    HTTP independiente.
    \item Las dos pestañas del navegador de una misma estación (cara
    Paciente y cara Recepción) se comunican entre sí mediante un canal de
    sincronización intra-navegador, sin cruzar ningún límite de nodo ni de
    proceso y sin tráfico de red.
    \item El sistema de archivos local de cada estación (\texttt{dataset/},
    \texttt{models/}, \texttt{voices/}) es leído y escrito únicamente por el
    proceso del backend de esa estación; no hay motor de base de datos ni
    almacenamiento compartido entre estaciones. La carpeta \texttt{logs/}
    (eventos de auditoría y consentimiento en \texttt{audit.jsonl}, reportes
    de traducción en \texttt{reports.jsonl}) también es escrita únicamente
    por ese proceso (RF-03, RF-04, RF-06).
    \item Dentro de cada estación, la cara Paciente usa WebSocket
    (\texttt{/ws/predict}) y HTTP REST contra su backend; la cara Recepción
    solo HTTP REST. El backend invoca a su propio Ollama local por HTTP
    (\texttt{POST /api/generate}).
\end{itemize}

% ============================================================================
\section{Vista de Escenarios (+1)}
\label{sec:vista-escenarios}

Esta vista integra las cuatro anteriores mediante dos escenarios de
ejecución concretos: \textbf{CU-01} (Comunicarse mediante Lengua de Señas
Argentina) y \textbf{CU-04} (Responder al Paciente). La Tabla~\ref{tab:arq-escenarios}
indica, para cada paso del escenario, qué elemento de cada vista participa.

\small
\begin{longtable}{>{\raggedright\arraybackslash}p{2.15cm} >{\raggedright\arraybackslash}p{2.35cm} >{\raggedright\arraybackslash}p{2.35cm} >{\raggedright\arraybackslash}p{2.15cm} >{\raggedright\arraybackslash}p{2.55cm} >{\raggedright\arraybackslash}p{1.75cm}}
    \toprule
    \textbf{Paso del escenario} & \textbf{Vista Lógica} & \textbf{Vista de Procesos} & \textbf{Vista de Desarrollo} & \textbf{Vista Física} & \textbf{RF / CU} \\
    \midrule
    \endhead
    Captura de landmarks &
    \texttt{useMediaPipe}, \texttt{lib/\allowbreak mediapipe.ts} &
    Bucle de captura (hilo principal del navegador) &
    \texttt{web/hooks}, \texttt{web/lib} &
    Pestaña Estación Paciente (nodo navegador) &
    RF-01.1, CU-01 \\

    Envío y predicción por WebSocket &
    \texttt{useSign\allowbreak Web\allowbreak Socket} $\to$ \texttt{app.py} &
    WS \texttt{/ws/predict} en el \textit{event loop}; inferencia en hilo propio &
    \texttt{web/hooks} $\to$ \texttt{app.py} (HTTP/WS) &
    Pestaña Paciente $\leftrightarrow$ proceso Uvicorn (\texttt{:8000}) &
    RF-01.2, RF-01.3, CU-01, CU-02 \\

    Cierre de oración y traducción &
    \texttt{useSentence\allowbreak Builder} $\to$ \texttt{app.py} &
    POST \texttt{/process\_sentence} en el \textit{threadpool} de Starlette; llamada a Ollama &
    \texttt{app.py} $\to$ \texttt{dataset/}, \texttt{models/} (indirectamente vía LSTM) &
    Proceso Uvicorn $\to$ proceso Ollama (\texttt{:11434}) &
    RF-01.4, RF-01.5, CU-02 \\

    Síntesis de voz &
    \texttt{app.py} (Piper embebido) &
    Llamada sincrónica a Piper dentro del \textit{threadpool} &
    \texttt{app.py} $\to$ \texttt{voices/} &
    Piper TTS embebido en el proceso Uvicorn &
    RF-01.6, CU-03 \\

    Publicación del mensaje del paciente &
    \texttt{useChat\allowbreak Channel} &
    Rama \texttt{par}: exhibición local + propagación por \texttt{Broadcast\allowbreak Channel} &
    \texttt{web/hooks} (sin dependencia de \texttt{app.py}) &
    Pestaña Paciente $\leftrightarrow$ Pestaña Recepción (intra-navegador) &
    RF-01.7, CU-01, CU-04 \\

    Respuesta del personal &
    \texttt{Reception\allowbreak Controls}, \texttt{useSpeech\allowbreak To\allowbreak Text}, \texttt{useChat\allowbreak Channel} &
    Evento \texttt{MESSAGE} propagado de vuelta a la Estación Paciente &
    \texttt{web/\allowbreak components}, \texttt{web/hooks} &
    Pestaña Recepción $\leftrightarrow$ Pestaña Paciente (intra-navegador) &
    RF-02.1, RF-02.2, CU-04, CU-05 \\
    \bottomrule
    \caption{Escenario CU-01 / CU-04 trazado a través de las cuatro vistas}
    \label{tab:arq-escenarios}
\end{longtable}

Este recorrido confirma que ningún paso del escenario requiere un nodo,
proceso o dependencia fuera de la PC local descripta en la Vista Física, lo
que sostiene la decisión de procesamiento local por privacidad (RNF-03).

% ============================================================================
\section{Decisiones arquitectónicas}
\label{sec:decisiones}

La Tabla~\ref{tab:arq-decisiones} resume las decisiones arquitectónicas más
relevantes de LSA Salud, su justificación y sus consecuencias sobre las
vistas anteriores.

\begin{longtable}{>{\raggedright\arraybackslash}p{4.25cm} >{\raggedright\arraybackslash}p{5.2cm} >{\raggedright\arraybackslash}p{5.2cm}}
    \toprule
    \textbf{Decisión} & \textbf{Justificación} & \textbf{Consecuencias} \\
    \midrule
    \endhead
    Procesamiento y persistencia 100\% locales (sin llamadas a servicios en
    la nube) &
    Los datos de salud del paciente no deben salir de la máquina local
    (RNF-03.1); evita depender de conectividad externa para el flujo crítico
    (RNF-T-02). &
    Ollama y el motor de síntesis de voz deben estar instalados de antemano
    en la máquina (ver ERS 2.7, Suposiciones); no hay escalado horizontal
    posible sin rediseño. \\

    Auditoría basada en archivos (\texttt{logs/audit.jsonl}) mediante
    funciones simples dentro de \texttt{app.py}
    (\texttt{\_log\_audit}/\texttt{\_append\_jsonl}), sin un módulo
    separado &
    RF-04.1 exige registrar logs de uso y métricas operativas
    anónimas para auditoría, sin agregar un servicio ni una base de datos
    nueva (coherente con la decisión de persistencia basada en archivos). &
    Cada \textit{handler} existente agrega una línea de invocación a
    \texttt{\_log\_audit}; el formato de línea (JSON Lines) permite anexar
    sin reescribir el archivo completo; la escritura comparte un
    \texttt{threading.Lock} con \texttt{/report}. \\

    Gateo de consentimiento dentro del propio \texttt{paciente/page}
    (\texttt{consentState} pasado como \texttt{enabled} al hook de
    landmarks), en vez de un componente de guardia separado &
    RF-03.1 exige bloquear la cámara hasta el consentimiento
    explícito; controlar el flag \texttt{enabled} del hook de captura desde
    la propia página evita introducir un componente y un árbol de render
    adicionales para un estado que ya vive en \texttt{PacientePage}. &
    \texttt{useMediaPipe} solo invoca \texttt{getUserMedia()} cuando
    \texttt{enabled === true}; ningún otro \texttt{hook} necesita conocer el
    estado de consentimiento. \\

    Streaming de frames al backend por WebSocket (\texttt{/ws/predict}) &
    Minimizar la latencia de ida y vuelta en el canal más sensible al
    rendimiento (RNF-01.1), evitando el costo de abrir una conexión HTTP por
    frame. &
    \texttt{useSignWebSocket} mantiene una única conexión persistente por
    sesión; el backend reconecta y descarta envíos obsoletos según
    RNF-T-01. \\

    MediaPipe como único backend de captura, extrayendo landmarks de manos,
    pose y rostro (\textit{blendshapes}) en un mismo paso por frame &
    Los \textit{blendshapes} faciales son necesarios para capturar los
    marcadores no manuales gramaticales de la LSA (preguntas, negación,
    intensidad), imprescindibles para la fidelidad clínica del vocabulario
    (RNF-05). &
    Un único \texttt{hook} de captura (\texttt{useMediaPipe}) y un único
    árbol de datos de entrenamiento (\texttt{dataset/custom}), sin
    variantes que mantener en paralelo. \\

    Restricción de CORS al origen de la red local del sanatorio en el
    backend (RNF-T-06) &
    Cada estación sirve únicamente a los dispositivos de su propia red
    local; no hay razón de negocio para aceptar conexiones desde otros
    orígenes, y restringirlas reduce la superficie de ataque. &
    El backend rechaza solicitudes con un origen no configurado; agregar un
    nuevo dispositivo a una estación requiere sumar su origen a la
    configuración local de esa estación. \\

    Persistencia basada en archivos (\texttt{.npy}, \texttt{.pth},
    \texttt{.onnx}) en lugar de una base de datos &
    El dominio no requiere consultas relacionales ni concurrencia de
    escritura alta: las muestras de entrenamiento y los modelos son
    artefactos versionables por archivo; simplifica el despliegue
    \textit{standalone} (ver Entregable 09). &
    Sin transacciones ni consultas indexadas; la consistencia de
    \texttt{dataset\_stats} depende de recorrer el árbol de carpetas en
    cada consulta. \\

    Backend de un solo archivo (\texttt{app.py}) para todos los endpoints &
    Simplicidad de despliegue en un entorno de consultorio, sin
    infraestructura de orquestación entre servicios. &
    Todos los endpoints comparten el mismo proceso, el mismo \textit{event
    loop} y el mismo \textit{threadpool}: una carga alta en un endpoint
    puede afectar la latencia de los demás (ver Vista de Procesos). \\

    Cachés LRU en memoria para traducción (\texttt{\_translation\_cache},
    200 entradas) y para audio sintetizado (\texttt{\_audio\_cache}, 50
    entradas) &
    Evitar repetir llamadas costosas al modelo de lenguaje o a la síntesis
    de voz ante oraciones ya traducidas o sintetizadas antes (RNF-01.1). &
    Las cachés son volátiles (en memoria del proceso): se pierden al
    reiniciar el backend; no hay invalidación explícita más allá del límite
    de tamaño (política LRU). \\
    \bottomrule
    \caption{Decisiones arquitectónicas de LSA Salud}
    \label{tab:arq-decisiones}
\end{longtable}

% ============================================================================
\section{Trazabilidad con atributos de calidad}
\label{sec:trazabilidad-rnf}

La Tabla~\ref{tab:arq-rnf} vincula cada requisito no funcional --- numerado
en la matriz de requisitos o definido por el equipo, ver \texttt{IDS.md}
Sección~4 --- con la vista o el mecanismo arquitectónico que lo atiende.

\begin{longtable}{>{\raggedright\arraybackslash}p{2.3cm} >{\raggedright\arraybackslash}p{2.8cm} >{\raggedright\arraybackslash}p{8.6cm}}
    \toprule
    \textbf{RNF} & \textbf{Atributo} & \textbf{Vista / mecanismo que lo atiende} \\
    \midrule
    \endhead
    RNF-T-01 & Fiabilidad &
    Vista de Procesos: reconexión automática del WebSocket cada 2 s,
    degradación sin interrupción si falla Ollama o Piper (ver
    sección~\ref{sec:vista-procesos}). \\

    RNF-T-02 & Disponibilidad &
    Vista Física: todos los procesos críticos corren en la PC local, sin
    dependencia de conectividad externa (sección~\ref{sec:vista-fisica}). \\

    RNF-03.1 & Seguridad / privacidad &
    Vista Física y decisión arquitectónica de procesamiento local
    (sección~\ref{sec:decisiones}): ningún dato de la conversación sale de
    la máquina; persistencia basada en archivos sin base de datos externa;
    diseño de auditoría y reporte (RF-04.1, RF-06.1) sin imágenes ni video. \\

    RNF-T-03 & Mantenibilidad &
    Vista Lógica: separación de responsabilidades en \texttt{hooks}
    independientes y módulo LSTM desacoplado del backend
    (sección~\ref{sec:vista-logica}). \\

    RNF-T-04 & Portabilidad &
    Vista Física: stack multiplataforma (Python/Uvicorn, Node.js/Next.js);
    probado en Windows (sección~\ref{sec:vista-fisica}). \\

    RNF-T-05 & Usabilidad / accesibilidad &
    Vista de Procesos: la ventana de exhibición previa a la reproducción de
    audio actúa como confirmación automática, sin botón manual
    (sección~\ref{sec:vista-procesos}). \\

    RNF-01.1 & Rendimiento &
    Vista de Procesos: separación event loop / \textit{threadpool}, y
    decisiones arquitectónicas de streaming por WebSocket y cachés LRU
    (secciones~\ref{sec:vista-procesos} y~\ref{sec:decisiones}). \\

    RNF-02.1 & Precisión &
    Vista de Desarrollo: la exactitud de validación se calcula en
    \texttt{train\_lstm.py} y se evalúa con \texttt{eval\_model.py},
    ambos fuera del proceso del backend (ver ERS, Sección~3.3). \\

    RNF-04.1 / RNF-07.1 & Hardware / estación fija &
    Vista Física: cada estación sincroniza sus dos caras y materializa la
    estación fija (sección~\ref{sec:vista-fisica}); el ensamblaje físico y
    la calibración se realizan durante la instalación de cada estación. \\

    RNF-05.1 / RNF-06.1 & Dataset clínico &
    Vista de Desarrollo: árbol único \texttt{dataset/custom}, con landmarks
    y \textit{blendshapes} en la misma muestra (sección~\ref{sec:vista-desarrollo});
    la validación lingüística por intérpretes de LSA es un procedimiento
    operativo del proyecto (ver ERS, Sección~2.6). \\

    RNF-T-06 & Seguridad de red &
    Vista Física y decisión arquitectónica de restricción de CORS: cada
    backend admite únicamente los orígenes de la red local de su propia
    estación (secciones~\ref{sec:vista-fisica} y~\ref{sec:decisiones}). \\
    \bottomrule
    \caption{Trazabilidad entre requisitos no funcionales y vistas arquitectónicas}
    \label{tab:arq-rnf}
\end{longtable}
